% DOCUMENT OPTIONS -------------------------------------------------------------
\makeatletter

\newcommand*{\installoption}[2][acmart]{
	\expandafter\newif\csname if#2\endcsname
	\@ifclasswith{#1}{#2}{
		\csname#2true\endcsname
	}{}
}
\newcommand*{\obeyoption}[2]{
	\csname if#2\endcsname\else\csname #1false\endcsname\fi
}
\newcommand*{\excludeoption}[2]{
	\csname if#2\endcsname\csname #1false\endcsname\fi
}
\makeatother

\installoption{highlight}
\installoption{appendix}
\installoption{final}
\installoption{draft}
\installoption{showframe}
\obeyoption{showframe}{draft}

% UTILS ------------------------------------------------------------------------
\usepackage{aliascnt}
\usepackage{xparse}
%\usepackage{xpatch}
\usepackage{etoolbox}
%\usepackage{xspace}
\usepackage{environ}
%\usepackage{ragged2e}
\usepackage{graphicx}

\def\deprecated#1{\PackageWarning{preamble}{Don't use ``#1''}{}}

% SPACING ----------------------------------------------------------------------

%\widowpenalty10000
%\clubpenalty10000

\def\proofltsskipamount{6pt}
\def\termltsskipamount{1\smallskipamount}
\def\envltsskipamount{1\smallskipamount}
\def\typerulesskipamount{1\smallskipamount}
%\def\rulescaling{.94}
%\newcommand\rulescalebox[2][\rulescaling]{\scalebox{#1}{$#2$}}
\providecommand\rulescalebox[2][]{#2} % no scaling

\makeatletter
\newlength{\@parindent}
\def\savepar{\setlength{\@parindent}{\parindent}}
\def\restorepar{\setlength{\parindent}{\@parindent}}
\makeatother

% PAGE&FONT --------------------------------------------------------------------
%\usepackage[T1]{fontenc}
\usepackage[utf8]{inputenc}
\usepackage{microtype}
\usepackage[normalem]{ulem}
\usepackage{slantsc}

\ifdefined\ifshowframe
  \ifshowframe
    \AtEndPreamble{%
      \usepackage{showframe}%
      \renewcommand*\ShowFrameColor{\color{black!50}}
      \renewcommand*\ShowFrameLinethickness{.5pt}
    }
  \fi
\fi

\newcommand\hmmax{0}
\newcommand\bmmax{0}

% COLOURS ----------------------------------------------------------------------
\usepackage{xcolor}

\colorlet{type}{blue!60!black}
\colorlet{process}{red!60!black}
\definecolor{light-gray}{gray}{0.8}

\colorlet{textcolor}{.}
\newcommand\usecolor[2]{\begingroup%
	\colorlet{uc@saved}{.}
	\color{#1}#2\color{uc@saved}
\endgroup}
\newcommand\cboxed[2]{\begingroup%
  \colorlet{cb@saved}{.}%
  \color{#1}%
  \boxed{\color{cb@saved}#2}%
\endgroup}

% TIKZ -------------------------------------------------------------------------
\usepackage{tikz}
\usepackage{hf-tikz} % highlighting boxes and annotations
\usetikzlibrary{calc}
\usetikzlibrary{arrows,matrix,positioning,scopes}

% BABEL ------------------------------------------------------------------------
\usepackage[british]{babel}
\usepackage[all]{foreign}
\renewcommand\foreignabbrfont{}
\renewcommand\foreignfullfont{}
%\defnotforeign[aka]{{aka}}
\defnotforeign[wrt]{{wrt}}
\defnotforeign[wl]{{wlog}}
%\defasforeign[etsimilia]{et similia}
%\defasforeign[perse]{per se}
%\defasforeign[perdefinitionem]{per definitionem}
%\defasforeign[perdef]{\xperiodafter{per def}}

%\hyphenation{}

% MATH -------------------------------------------------------------------------
\usepackage{amsmath}
\usepackage{amsthm}
\usepackage{amsfonts}
% \usepackage{amssymb}
\usepackage{stmaryrd}
\SetSymbolFont{stmry}{bold}{U}{stmry}{m}{n}
\usepackage{bm}
\usepackage{mathtools}
\usepackage{centernot}
\usepackage{scalerel}

%\makeatletter
%\newcommand{\leqnomode}{\tagsleft@true}
%\newcommand{\reqnomode}{\tagsleft@false}
%\makeatother

\makeatletter
\newcommand{\pushright}[1]{\ifmeasuring@#1\else\omit\hfill$\displaystyle#1$\fi\ignorespaces}
\newcommand{\pushleft}[1]{\ifmeasuring@#1\else\omit$\displaystyle#1$\hfill\fi\ignorespaces}
\newcommand{\specialcell}[1]{\ifmeasuring@#1\else\omit$\displaystyle#1$\ignorespaces\fi}
\makeatother

\allowdisplaybreaks

%% extensible rightarrowtriangle
\makeatletter
\@ifundefined{minaw@}{\newdimen\minaw@}{}
\def\rightarrowtrianglefill@#1{\m@th\setboxz@h{$#1\relbar$}\ht\z@\z@
  $#1\copy\z@\mkern-6mu\cleaders
  \hbox{$#1\mkern-2mu\box\z@\mkern-2mu$}\hfill
  \mkern-6mu\mathord{\smash\rightarrowtriangle\vphantom{\rightarrow}}$}
\newcommand{\xrightarrowtriangle}[2][]{%
  \mathrel{\mathop{%
    \setbox\z@\vbox{\m@th
      \hbox{$\scriptstyle\;{#1}\;\;\,$}%
      \hbox{$\m@th\scriptstyle\;{#2}\;\;\,$}%
    }%
    \hbox to\ifdim\wd\z@>\minaw@\wd\z@\else\minaw@\fi{%
      \rightarrowtrianglefill@\displaystyle}}%
  \limits^{#2}\@ifnotempty{#1}{_{#1}}}%
}
% FIX \xRightarrow text padding
\renewcommand{\xRightarrow}[2][]{\ext@arrow 0359\Rightarrowfill@{#1}{#2}}
\makeatother

% load arrows with stealth heads from mathabx
%\DeclareFontFamily{U}{matha}{\hyphenchar\font45}
%\DeclareFontShape{U}{matha}{m}{n}{
%     <5> <6> <7> <8> <9> <10> gen * matha
%     <10.95> matha10 <12> <14.4> <17.28> <20.74> <24.88> matha12
%     }{}
%\DeclareSymbolFont{matha}{U}{matha}{m}{n}
%\DeclareMathSymbol{\varleftarrow}{3}{matha}{"D0}
%\DeclareMathSymbol{\varrightarrow}{3}{matha}{"D1}
%\DeclareMathSymbol{\varleftrightarrow}{3}{matha}{"D8}

\def\lclens{{[\llap(}}
\def\rclens{{\rlap)]}}

%% load \lBrace, \rBrace from stix
%\DeclareFontEncoding{LS2}{}{\noaccents@}
%\DeclareFontSubstitution{LS2}{stix}{m}{n}
%\DeclareSymbolFont{stix@largesymbols}{LS2}{stixex}{m}{n}
%\SetSymbolFont{stix@largesymbols}{bold}{LS2}{stixex}{b}{n}
%\DeclareMathDelimiter{\lBrace}{\mathopen} {stix@largesymbols}{"E8}%
%                                          {stix@largesymbols}{"0E}
%\DeclareMathDelimiter{\rBrace}{\mathclose}{stix@largesymbols}{"E9}%
%                                          {stix@largesymbols}{"0F}

\makeatletter
\newcommand{\raisemath}[1]{\mathpalette{\raisem@th{#1}}}
\newcommand{\raisem@th}[3]{\raisebox{#1}{$#2#3$}}
\makeatother



% TABULARS ---------------------------------------------------------------------
\usepackage{array}
%\usepackage{xtab}
%\xentrystretch{-0.2}
\usepackage{longtable}
\usepackage{booktabs}
\usepackage{multicol}
%\usepackage{colortbl}
%\usepackage{arydshln}

\makeatletter
\def\crulefill#1{\leavevmode\leaders\hrule\@height#1 height 0.7ex depth \dimexpr0.4pt-0.7ex\hfill\kern0pt}
\makeatother

%\def\cdashfill{\leavevmode\cleaders\hbox{--}\hfill\kern0pt}

\def\headertext#1{
	\vspace{2pt}
	\crulefill{\lightrulewidth}\ \textsf{\small \emph{#1}}\ \crulefill{\lightrulewidth}
  \\[-2.7ex]%
% 	\rule[-1ex]{\textwidth}{\heavyrulewidth}%
% 	\vspace{-1.6\smallskipamount}%
}

\def\header{
 	\vspace{-1.6\smallskipamount}%
 	\rule[1\abovedisplayskip]{\textwidth}{\heavyrulewidth}%
}

\newcommand*{\footer}[1][1.6\belowdisplayskip]{%
	\rule[#1]{\textwidth}{\lightrulewidth}%
	\vspace{-1.6\smallskipamount}%
}

% FLOATS -----------------------------------------------------------------------
\usepackage{float}
\usepackage{placeins}
\usepackage{wrapfig}

% ENUMS ------------------------------------------------------------------------
\usepackage[inline]{enumitem}

% CROSS REFS -------------------------------------------------------------------
\usepackage{nameref}

\usepackage[capitalise,nameinlink,noabbrev]{cleveref}
 % extend cref with custom labels
\makeatletter
% #1 number (used sorting and intervals by cref); 
% #2 type;  (used by cref for name and format)
% #3 label; 
% #4 tag
\newcommand{\customlabel}[4][0]{%
	\protected@write\@auxout{}{\string\newlabel{#3}{{#4}{\thepage}{#4}{#3}{}}}%
	\protected@write\@auxout{}{\string\newlabel{#3@cref}{{[#2][#1][#1]#4}{\thepage}{}{}{}}}%
}
\newcommand{\phantomlabel}[4][0]{%
	\customlabel[#1]{#2}{#3}{#4}%
	\hypertarget{#3}{}%
}
\newcommand{\labelalias}[4][0]{
	\customlabel[#1]{#2}{#3}{#4}%
}
\newcommand{\crefv}[1]{%
	\begingroup\@cref@compressfalse\@cref@sortfalse\cref{#1}\endgroup%
}
\newcommand{\Crefv}[1]{%
	\begingroup\@cref@compressfalse\@cref@sortfalse\Cref{#1}\endgroup%
}
\newcommand{\crefabbrev}[1]{%
	\begingroup\@cref@abbrevtrue\cref{#1}\endgroup%
}
\makeatother

\AtEndPreamble{\hypersetup{
%	hidelinks,
	final,
}}

% RULES & DERIVATIONS ----------------------------------------------------------

\usepackage{proof-dashed}

% internal counter for ordering rules in \cref by order of definition
\newcounter{rule}

% rule names
\newcommand{\rname}[1]{\ensuremath{\textsc{\MakeLowercase{#1}}}}
% rule labels
% #1 tag; #2 label
\newcommand{\rlabel}[2]{{%
	\refstepcounter{rule}
	\customlabel[\therule]{infrule}{#2}{#1}%
	\hypertarget{#2}{#1}}}
% cref support
\crefname{infrule}{rule}{rules}
\Crefname{infrule}{Rule}{Rules}
%\creflabelformat{infrule}{#2\textup{(#1)}#3}

% marking old rules
\newcommand{\rrelabel}[2]{\rlabel{\ref*{#1}}{#2}}
% forward a label
% #1 target label (existing)
% #2 alias label  (new)
\newcommand{\rlabelfwd}[2]{
	\refstepcounter{rule}
	\customlabel[\therule]{infrule}{#2}{\ref*{#1}}%
}

% example:
% \infer
%   [\rlabel{\rname{Awesome!}}{rule:awesome}]
%   {conclusion}
%   {premises}
% \cref{rule:awesome} -> Rule Awesome!

% THEOREMS ET SIMILIA ----------------------------------------------------------

% restatable theorems
\usepackage{thmtools}
\usepackage{thm-restate}

\AtEndPreamble{
	\theoremstyle{acmplain}
	%\newtheorem{theorem}{Theorem}
	%\newtheorem{corollary}[theorem]{Corollary}
	%\newtheorem{proposition}[theorem]{Proposition}
	%\newtheorem{lemma}[theorem]{Lemma}
  \newtheorem{exercise}[theorem]{Exercise}
	\newtheorem{fact}[theorem]{Fact}
	\theoremstyle{acmdefinition}
	%\newtheorem{definition}[theorem]{Definition}
	%\newtheorem{assumption}[definition]{Assumption}
	%\newtheorem{convention}[definition]{Convention}
	%\newtheorem*{notation}{Notation}
	\newtheorem{remark}[theorem]{Remark}
}

\crefname{fact}{Fact}{Facts}
\Crefname{fact}{Fact}{Facts}

\crefname{proof}{proof of}{proofs of}
\Crefname{proof}{Proof of}{Proofs of}

% #1 theorem; #2 label
\newcommand{\prooflabel}[2]{{%
	\customlabel{proof}{#2}{\cref*{#1}}%
	\hypertarget{#2}{}}}

% Biblio -----------------------------------------------------------------------
\bibliographystyle{ACM-Reference-Format}
\citestyle{acmauthoryear}

% REVISION ---------------------------------------------------------------------
\usepackage[
	disable,     % hide all todos
%	obeyDraft,   % show todos only in draft mode
	obeyFinal,   % hide todos if in final mode
	textsize=tiny,
	colorinlistoftodos,
	prependcaption
]{todonotes}
%\newcommand{\unsure}[2][]{\todo[linecolor=red,backgroundcolor=red!25,bordercolor=red,#1]{#2}}
%\newcommand{\change}[2][]{\todo[linecolor=blue,backgroundcolor=blue!25,bordercolor=blue,#1]{#2}}
%\newcommand{\info}[2][]{\todo[linecolor=cyan,backgroundcolor=cyan!25,bordercolor=cyan,#1]{#2}}
%\newcommand{\improvement}[2][]{\todo[linecolor=green,backgroundcolor=green!25,bordercolor=green,#1]{#2}}
\usepackage{verbatim}

% MACROS -----------------------------------------------------------------------

%- acronyms --------------------------------------
\makeatletter
\def\new@acronym#1#2{%
  \expandafter\def\csname#1@acr\endcsname{\text{\upshape #2}\xspace}%
  \expandafter\def\csname#1\endcsname{{\upshape #2}\xspace}}
\def\new@acronymexp#1#2{\expandafter\def\csname#1@exp\endcsname{{#2}}}
\NewDocumentCommand{\newacronym}{o m o}{%
	\IfNoValueTF{#1}{\def\n{#2}}{\def\n{#1}}
  \new@acronym{\n}{#2}
  \IfNoValueF{#3}{\new@acronymexp{\n}{#3}}
}
\def\acronym#1{\csname#1@acr\endcsname}
\def\expandacronym#1{\csname#1@exp\endcsname}
\makeatother

\newacronym{CP}[Classical Processes]
\newacronym{MCP}[Multiplicative Classical Processes]
\newacronym[CPZero]{CP\textsubscript{0}}[Non-Polymorphic Classical Processes]
\newacronym{CHOP}[Classical Higher-Order Processes]
\newacronym{HCP}[Hypersequent Classical Processes]
\newacronym{DHCP}[Hypersequent Classical Processes with Delayed Actions]
\newacronym[HOpi]{HO$\pi$}[Higher-Order $\pi$-calculus]
\newacronym{CLL}[Classical Linear Logic]
\newacronym{MCLL}[Multiplicative Classical Linear Logic]
\newacronym[CLLZero]{CLL\textsubscript{0}}[Quantifier-free Classical Linear Logic]
\newacronym[CLLZeroOne]{CLL\textsubscript{01}}[Classical Linear Logic with first-order quantifiers]
\newacronym{ILL}[Intuitionistic Linear Logic]
\newacronym[lts]{lts}[labelled transition system]
\newacronym[pill]{$\pi$LL\xspace}
\newacronym[hopill]{HO$\pi$LL\xspace}[Higher-Order \pill]
\newacronym[polypill]{Poly$\pi$LL\xspace}[Polymorphic \pill]
\newacronym[pimll]{$\pi$MLL\xspace}[Multiplicative \pill]
%\newacronym[pimell]{$\pi$MELL\xspace}
%\newacronym[pimall]{$\pi$MALL\xspace}
\newacronym[repill]{Re$\pi$LL\xspace}[\pill with Recursion]

%- operators -------------------------------------
\DeclareMathOperator{\dom}{dom}
\DeclareMathOperator{\cod}{cod}
\DeclareMathOperator{\img}{img} % image
\DeclareMathOperator{\pre}{pre} % preimage

\DeclareMathOperator{\prfhyp}{hyp}
\DeclareMathOperator{\prfproc}{proc}
\DeclareMathOperator{\cn}{n} % channel names (free + bound)
\DeclareMathOperator{\fn}{fn} % free names
\DeclareMathOperator{\bn}{bn} % bound names
\DeclareMathOperator{\pv}{p} % process variables
\DeclareMathOperator{\fp}{fp} % free process variables
\DeclareMathOperator{\bp}{bp} % bound process variables
\DeclareMathOperator{\ft}{ft} % free type variables
\DeclareMathOperator{\bt}{bt} % bound type variables
\DeclareMathOperator{\fs}{f} % free symbols (names+vars)
\DeclareMathOperator{\bs}{b} % bound symbols
\DeclareMathOperator{\is}{i} % introduced symbols


\makeatletter
\DeclareMathOperator{\bigtensor}{{\textstyle\bigotimes}}
\DeclareMathOperator{\bigparr}{\scalerel*{\parr}{\bigtensor}}
\makeatother

\makeatletter
\newcommand{\enc}{\@ifstar{\enc@}{\enc@@}}
\def\enc@#1{\left\llbracket \mkern-7mu\array{c}#1\endarray\mkern-7mu \right\rrbracket}
\def\enc@@#1{\llbracket#1\rrbracket}
\makeatother

\DeclareMathOperator{\disenop}{dis}  % disentanglement
\DeclareMathOperator{\packop}{pack}  % packing
\DeclareMathOperator{\unpackop}{unpack}

\newcommand{\disen}[1]{\disenop(#1)}
\newcommand{\pack}[3][]{\packop^{#1}_{#2}(#3)}
\newcommand{\unpack}[3][]{\unpackop^{#1}_{#2}(#3)}

\def\ppcount#1{\left\|#1\right\|} % num of parallel components 
%\def\ppdrop#1#2{\mathrm{drop}_{#1}(#2)}
%\def\ppkeep#1#2{\mathrm{keep}_{#1}(#2)}
\def\ppdrop#1#2{\left\lceil#2\right\rceil_{#1}}
\def\ppkeep#1#2{\left\lfloor#2\right\rfloor^{#1}}

%- relations -------------------------------------

% calligraphic letters for relations
\def\crel#1{\mathrel{\mathcal{#1}}}

% first and second projections
\def\fst{\mathop{\mathit{fst}}}
\def\snd{\mathop{\mathit{snd}}}

% disjointness
\def\disjoint{\mathrel{\#}}

% definitions
\let\defeq\triangleq
\edef\defiff{\stackrel{\triangle}{\iff}}

% alpha eq
\edef\aleq{\mathbin{=_{\alpha}}}
% structural congruence
\let\congr\equiv
% type equivalence
\def\tyeq{\mathrel{\dashv\vdash}}
% bisimulation, strong bis, barbed cong, forward congruence
\let\bis\approx
\let\sbis\sim
\let\bcong\approxeq
% denotational equivalence
\let\deneq\bumpeq
% strong similarity
\def\ssimby{\lesssim}
\def\ssims{\gtrsim}
% (weak) similarity
\def\simby{\lessapprox}
\def\sims{\gtrapprox}

% reductions
\let\reducesto\rightarrow
\def\breducesto{\mathbin{{\rightarrow}_\beta}}
\def\kreducesto{\mathbin{{\rightarrow}_\kappa}}

% lts arrows
\makeatletter
\def\lto#1{\xrightarrow{\lbox{#1}}} % labelled transitions
\def\slto#1{\xRightarrow{\lbox{#1}}} % saturated ones
\def\notlto#1{\centernot{\xrightarrow{\lbox{#1}}}}
\def\notslto#1{\centernot{\xRightarrow{\lbox{#1}}}}
% minimum width for labels
\newcommand\lbox[2][.8em]{
	\sbox0{$\scriptstyle #2$}
	\ifdim \wd0 < #1%
		\mathmakebox[#1]{\usebox0}%
	\else%
		\mathmakebox[\wd0]{\usebox0}%
	\fi%
}

% labelled transitions between proofs.
% * vertical layout; #1 source; #2 label; #3 target.
\newcommand{\plto}{%
	\@ifstar{\plto@v\downarrow}{\plto@h\lto}%
}
% saturated labelled transitions between proofs.
\newcommand{\pslto}{%
	\@ifstar{\plto@v\Downarrow}{\plto@h\slto}%
}
\def\pboxcolor{black!10}
\def\pbox#1{\cboxed{\pboxcolor}{\mspace{-10mu}\array{c}#1\endarray\mspace{-10mu}}}
\let\plto@b\pbox
% #1 arrow; #2 source; #3 label; #4 target.
\def\plto@h#1#2#3#4{%
	\plto@b{#2}%
	#1{#3}%
	\plto@b{#4}%
}
\def\plto@v#1#2#3#4{%
	\array[b]{l}%
		\plto@b{#2}%
		\\%
		\quad%
		\cboxed{white}{\scriptstyle\left#1\cboxed{white}{\scriptstyle#3}\right.}\\%
		\plto@b{#4}%
	\endarray%
}
\newcommand{\notplto}[2]{%
	\plto@b{#1} \notlto{#2}%
}
\def\pvlto#1{\quad\cboxed{white}{\scriptstyle\left\plto@v\downarrow\cboxed{white}{\scriptstyle#1}\right.}}%
\def\pvlto#1{\quad\cboxed{white}{\scriptstyle\left\plto@v\Downarrow\cboxed{white}{\scriptstyle#1}\right.}}%
\makeatother

\newenvironment{drules}
  {\highlight\small\spreadlines{\proofltsskipamount}\csname gather*\endcsname}
  {\csname endgather*\endcsname\endspreadlines\endhighlight}

% TODO: add parenthesis like in eqref
\crefname{ltstrace}{execution}{executions}
\Crefname{ltstrace}{Execution}{Executions}
\crefalias{equation}{ltstrace}

%- preds -----------------------------------------

\def\terminates{{\downharpoonright}}
\def\diverges{{\upharpoonright}}
\def\mayterminate{{\Downharpoonright}}
\def\maydiverge{{\Upharpoonright}}

%- languages -------------------------------------

% highlighting for processes and types/propositions
% 	enabled in environment highlight or \hl
% 	globally governed by the option highlight
\def\phl#1{#1}% process highlight
\def\thl#1{#1}% type highlight
\def\nohl#1{#1}% no highlight
\newcommand{\hl}[1]{\highlight#1\endhighlight}
\ifhighlight
\newenvironment{highlight}
	{\ignorespaces%
	 \def\phl##1{\usecolor{process}{##1}}%
	 \def\thl##1{\usecolor{type}{##1}}%
	 \def\nohl##1{\usecolor{textcolor}{##1}}}%
	{\def\phl##1{##1}%
   \def\thl##1{##1}%
   \def\nohl##1{##1}%
	 \ignorespacesafterend}
\else
\newenvironment{highlight}{\ignorespaces}{\ignorespacesafterend}
\fi

\newenvironment{bnftable}[3][\phl]
{\highlight%
\setlength\LTleft{0pt}%
\setlength\LTright{0pt}%
\sbox0{$#2$}%
\sbox1{$#3$}%
\longtable{%
  >{\raggedleft$}m{\wd0}<{$}%
  >{\raggedright$#1\begingroup}m{\wd1}<{\endgroup$}%
  >{\it}l<{}}}
{\endlongtable\endhighlight}

\newenvironment{bnftable2col}[4][\phl]
{\highlight%
\setlength\LTleft{0pt}%
\setlength\LTright{0pt}%
\sbox0{$#2$}%
\sbox1{$#3$}%
\sbox2{$#4$}%
\longtable{%
  >{\raggedleft$}m{\wd0}<{$}%
  >{\raggedright$#1\begingroup}m{\wd1}<{\endgroup$}%
  >{\it}l<{}
	>{$}r<{$}%
  >{\raggedright$#1\begingroup}m{\wd2}<{\endgroup$}%
  >{\it}l<{}}}
{\endlongtable\endhighlight}

\newenvironment{bnftable2}[4][\phl]
{\highlight%
\setlength\LTleft{0pt}%
\setlength\LTright{0pt}%
\sbox0{$#2$}%
\sbox1{$#3$}%
\sbox2{#4}%
\longtable{%
  >{\raggedleft$}m{\wd0}<{$}%
  >{\raggedright$#1\begingroup}m{\wd1}<{\endgroup$}%
  >{\it}l<{}@{\extracolsep{\fill}}%
  >{\hfill\it}m{\wd2}<{}}}
{\endlongtable\endhighlight}
%	{\vskip\abovedisplayskip\rmfamily\flushleft\highlight%
%	  \sbox0{$#2$}%
%	  \sbox1{$#3$}%
% 	  \sbox2{#4}%
%%	  \xentrystretch{-0.1}%
%	  \begin{xtabular*}{\textwidth}{%
%	      >{\raggedleft$}m{\wd0}<{$}%
%	      >{\raggedright$#1\begingroup}m{\wd1}<{\endgroup$}%
%	      >{\it}l<{}@{\extracolsep{\fill}}
%	      >{\hfill\it}m{\wd2}<{}}}
%	{\end{xtabular*}\endhighlight\endflushleft\vskip\belowdisplayskip}

%\newenvironment{bnftable}[3][\phl]
%	{\vskip\abovedisplayskip\rmfamily\flushleft\highlight%
%	  \sbox0{$#2$}%
%	  \sbox1{$#3$}%
%%	  \xentrystretch{-0.1}%
%	  \begin{xtabular*}{\textwidth}{%
%	      >{\raggedleft$}m{\wd0}<{$}%
%	      >{\raggedright$#1\begingroup}m{\wd1}<{\endgroup$}%
%	      >{\it}l<{}}}
%	{\end{xtabular*}\endhighlight\endflushleft\vskip\belowdisplayskip}
%
%\newenvironment{bnftable2}[4][\phl]
%	{\vskip\abovedisplayskip\rmfamily\flushleft\highlight%
%	  \sbox0{$#2$}%
%	  \sbox1{$#3$}%
% 	  \sbox2{#4}%
%%	  \xentrystretch{-0.1}%
%	  \begin{xtabular*}{\textwidth}{%
%	      >{\raggedleft$}m{\wd0}<{$}%
%	      >{\raggedright$#1\begingroup}m{\wd1}<{\endgroup$}%
%	      >{\it}l<{}@{\extracolsep{\fill}}
%	      >{\hfill\it}m{\wd2}<{}}}
%	{\end{xtabular*}\endhighlight\endflushleft\vskip\belowdisplayskip}

%- types and environments ------------------------

\newcommand*{\bang}{\mathord{!}}
\newcommand*{\query}{\mathord{?}}
\newcommand*{\tensor}{\otimes}
\newcommand*{\parr}{\mathbin{\bindnasrepma}}
\newcommand*{\with}{\mathbin{\binampersand}}
\newcommand*{\one}{1}
\newcommand*{\zero}{0}
\newcommand*{\dual}[1]{#1^{\bot}}
\newcommand{\provide}[1]{{[#1]}}
\newcommand{\assume}[1]{{\langle #1 \rangle}}

\newcommand*{\cht}[2]{\phl{#1}\nohl{\colon}\mspace{-2mu}\thl{#2}} % channel type #1 : #2
\newcommand*{\pvt}[2]{\phl{#1}\nohl{\colon}\mspace{-2mu}\thl{#2}} % process variable type #1 : #2
\newcommand*{\emptyseq}{\bullet}
\newcommand*{\emptyhyp}{\varnothing}
\newcommand*{\emptyho}{\cdot}
\newcommand*{\hyp}[1]{\mathcal{#1}}
%\newcommand*{\ecx}[1]{\mathcal{#1}}

\let\seq\vdash
\newcommand*{\der}[1]{\ensuremath{\mathcal{#1}}}
\let\prf\der
\def\ccolon{\mathrel{:\mspace{-1mu}:}}
\newcommand{\judge}[4][]{{{}\thl{#2} \nohl{\seq_{\thl{#1}}} \phl{#3} \nohl{\ccolon} \thl{#4}}}
\newcommand{\cpjudge}[2]{{{}\nohl{\seq_{\text{\tiny\CP}}} \phl{#1} \nohl{\ccolon} \thl{#2}}}

\def\tw{\mathop{\mathrm{u}}}
\def\twstrip#1{\llparenthesis#1\rrparenthesis}
\newcommand{\twjudge}[4]{{\thl{#2} \nohl{\seq^{#1}} \phl{#3} \nohl{\ccolon} \thl{#4}}}

\def\twrname#1{{#1\textsuperscript{\tiny$\tw$}}} % \rname for potential
\def\cprname#1{{#1\textsuperscript{\tiny\CP}}} % \rname for CP

\def\DerSet{\mathit{Der}} % set of all typing derivations 
\def\EnvSet{\mathit{Env}} % set of all typing environments
\def\env{\mathit{env}} % project derivations to environments

\newcommand{\hoenv}[2]{{\langle\thl{#1};\thl{#2}\rangle}}

%- processes -------------------------------------
\newcommand*{\prc}[1]{\mathit{#1}}
%- messages
\newcommand*{\msg}[1]{\ensuremath{\text{\upshape\scshape{\MakeLowercase{#1}}}}}
\newcommand*{\minl} {\msg{l}}
\newcommand*{\minr} {\msg{r}}
\newcommand*{\muse} {\msg{use}}
\newcommand*{\mdisp}{\msg{disp}}
\newcommand*{\mdup} {\msg{dup}}
%- prefixes
\newcommand*{\asend}[2]{{#1[#2]}}
\newcommand*{\arecv}[2]{{#1(#2)}}
\newcommand*{\aclose}[1]{\asend{#1}{}}
\newcommand*{\await}[1] {\arecv{#1}{}}
\newcommand*{\ainl}[1]{\asend{#1}{\minl}}
\newcommand*{\ainr}[1]{\asend{#1}{\minr}}
\let\asendtype\asend
\let\arecvtype\arecv
\newcommand*{\aclientuse}[1] {\asend{#1}{\muse}}
\newcommand*{\aclientdisp}[1]{\asend{#1}{\mdisp}}
\newcommand*{\aclientdup}[2]{\asend{#1}{\mdup}(#2)}
\newcommand*{\forwardop}{{\leftrightarrow}}
\newcommand*{\aforward}[3][]{{#2 \forwardop^{#1} #3}}
\let\asendho\asend
\let\arecvho\arecv
\newcommand{\asendfree}[2]{#1 \langle #2 \rangle}
%- processes
\newcommand{\nil}{\mathbf{0}}
\let\forward\aforward
\newcommand*{\pp}{\mathbin{\boldsymbol{|}}}
\newcommand*{\resop}{\mathop{\boldsymbol{\nu}}}
\newcommand*{\res}[2]{{\boldsymbol{\nu}{#1}\mspace{2mu}}{#2}}
\newcommand*{\prefixed}[2]{{#1{.}#2}}
\newcommand*{\send}[3]{\prefixed{\asend{#1}{#2}}{#3}}
\newcommand*{\recv}[3]{\prefixed{\arecv{#1}{#2}}{#3}}
\newcommand*{\close}[2]{\prefixed{\aclose{#1}}{#2}}
\newcommand*{\wait}[2]{\prefixed{\await{#1}}{#2}}
\newcommand*{\inl}[2]{\prefixed{\ainl{#1}}{#2}}
\newcommand*{\inr}[2]{\prefixed{\ainr{#1}}{#2}}
\def\achoice#1{{#1}.\mathsf{case}}
\NewDocumentCommand{\choice}{s m m m}{
	\IfBooleanTF{#1}
	{\achoice{#2}\mspace{-2mu}\left\{\mspace{-10mu}\array{rl}%
		\minl\colon&\mspace{-15mu}#3\\%
		\minr\colon&\mspace{-15mu}#4%
	\endarray\mspace{-8mu}\right\}}%
	{\achoice{#2}\mspace{-2mu}\left\{%
		\minl\colon\mspace{-4mu}#3,%
		\minr\colon\mspace{-4mu}#4\right\}}}
\NewDocumentCommand{\choicebit}{s m m m}{
	\IfBooleanTF{#1}
	{\achoice{#2}\left\{\mspace{-10mu}\array{rl}%
		0\mapsto&\mspace{-15mu}#3\\%
		1\mapsto&\mspace{-15mu}#4%
	\endarray\mspace{-8mu}\right\}}%
	{\achoice{#2}\left\{%
		0\mapsto\mspace{-4mu}#3,%
1\mapsto\mspace{-4mu}#4\right\}}}
\newcommand*{\sendtype}[3]{\prefixed{\asendtype{#1}{#2}}{#3}}
\newcommand*{\recvtype}[3]{\prefixed{\arecvtype{#1}{#2}}{#3}}
\newcommand*{\clientdisp}[2]{\prefixed{\aclientdisp{#1}}{#2}}
\newcommand*{\clientuse}[2]{\prefixed{\aclientuse{#1}}{#2}}
\newcommand*{\clientdup}[3]{\prefixed{\aclientdup{#1}{#2}}{#3}}
% \newcommand*{\serverOld}[2]{{\fboxsep=0pt\colorbox{yellow}{\ensuremath{\serverNoDep{#1}{#2}}}}}
\newcommand*{\serverNoDep}[2]{{\mathsf{!}#1}.\mspace{-2mu}\left\{{#2}\right\}}
\newcommand*{\server}[3]{{\mathsf{!}#1}{\langle{#2}\rangle}.\mspace{-2mu}\left\{{#3}\right\}}

\newcommand*{\chopop}[2]{{[#1\,{\coloneqq}\,#2]}}
\newcommand*{\chop}[3]{{{#3}\chopop{#1}{#2}}}
\newcommand*{\sendho}[3]{\prefixed{\asendho{#1}{#2}}{#3}}
\newcommand*{\recvho}[3]{\prefixed{\arecvho{#1}{#2}}{#3}}
\newcommand*{\invoke}[3][]{{#2} \langle #3 \rangle}
\newcommand*{\abstr}[2]{(#1)#2}
\makeatletter
\newcommand{\pabstr}{\@ifstar{\pabstr@}{\pabstr@@}}
\newcommand*{\pabstr@}[2]{(#1)\cboxed{black!20}{%
  \let\@ots\inferTabSkip%
  \def\inferTabSkip{\hspace{1.4ex}}%
  \scriptstyle#2}}
\newcommand*{\pabstr@@}[4]{\pabstr@{#1}{\judge{#2}{#3}{#4#1}}}
\makeatother
\newcommand{\sendfree}[3]{\prefixed{\asendfree{#1}{#2}}{#3}}

\NewDocumentCommand{\hole}{o}{\IfNoValueTF{#1}{\Box}{\boxed{#1}}}
\let\ctx\mathbb

\def\ProcSet{\mathit{Proc}} % set of all processes
\def\proc{\mathit{proc}} % project derivations to processes

%- labels ----------------------------------------

\def\ActSet{\mathit{Act}} % set of action labels
\def\LblSet{\mathit{Lbl}} % set of all labels

%\newcommand{\co}[1]{\overline{#1}}

\newcommand*{\lsync}[2]{{#1{\pp}#2}}
\newcommand*{\lforward}[2]{{\aforward{#1}{#2}}}
\newcommand*{\lsend}[2]{{\asend{#1}{#2}}}
\newcommand*{\lrecv}[2]{{\arecv{#1}{#2}}}
\newcommand*{\lclose}[1]{{\aclose{#1}}}
\newcommand*{\lwait}[1]{{\await{#1}}}
\newcommand*{\lsendtype}[2]{{\asendtype{#1}{#2}}}
\newcommand*{\lrecvtype}[2]{{\arecvtype{#1}{#2}}}
\newcommand*{\linl}[1]   {{\ainl{#1}}}
\newcommand*{\lcoinl}[1] {{\lrecv{#1}{\minl}}}
\newcommand*{\linr}[1]   {{\ainr{#1}}}
\newcommand*{\lcoinr}[1] {{\lrecv{#1}{\minr}}}
\newcommand*{\luse}[1]   {{\aclientuse{#1}}}
\newcommand*{\lcouse}[1] {{\lrecv{#1}{\muse}}}
\newcommand*{\ldup}[1]   {{\lsend{#1}{\mdup}}}
\newcommand*{\lcodup}[1] {{\lrecv{#1}{\mdup}}}
\newcommand*{\ldisp}[1]  {{\lsend{#1}{\mdisp}}}
\newcommand*{\lcodisp}[1]{{\lrecv{#1}{\mdisp}}}
\newcommand*{\lsendho}[2]{{\asendho{#1}{#2}}}
\newcommand*{\lrecvho}[2]{{\arecvho{#1}{#2}}}

\newcommand*{\lsubst}[3][]{\chopop{{#2}}{#3}}
\newcommand*{\lchop}[2]{{#1\nohl{\seq\ccolon}#2}}

\def\plbl#1#2{#1} % used in HCP papers (deprecated)
%\let\alignatthis\relax % a little hack inject & in \plbl
%\newcommand{\plbl}[2]{\phl{#1}\colon\mspace{-2mu}\thl{#2}}

% categories and functors
\def\Set{\mathsf{Set}}
\def\pw{\mathcal{P}}
\def\lb{\mathcal{L}}

\newcommand{\newtext}[1]{{\color{blue} #1}}